\cotefile{34}
\scanpage{100}
\mspage{15}
% archive sheet number: 99
% source lines 1--4
sans supposer $U$ quasi-projectif sur $K$, mais seulement de t.f. et séparé. D'autre part, les réductions \deleted{faites à la fin de} \inserted{dans} I nous permettent de conclure :

% source lines 5--9
\uline{Corollaire} 3.7 Sous les conditions (i) (ii) de 3.5, pour tt schéma \inserted{$U$} de type \inserted{fini} sur $K$, $H^i(U_{\bar K},\unsure{\mathbb Z_\ell|Z_\ell})$ est de t.f. sur $\mathbb Z_\ell$ et l'action du groupe d'inertie $\bar I$ sur \deleted{\unsure{H}} lui est \unsure{qsi|quasi} unipotente. Même \unsure{résultats|résultat} pour $H^i_!(U_{\bar K},\mathbb Z_\ell)$, quand $U$ est \deleted{\unsure{de t.f.}} \inserted{séparé} sur $K$.

% source lines 10--23, including the interlinear revisions
\uline{Remarques} 3.8. \deleted{\unsure{La démonstration que nous donnons}} \inserted{\unsure{a)|a°)} Pour l'instant, 3.5 \deleted{\unsure{n'est applicable}} \inserted{et 3.7,}} \unsure{sont applicables} tels quels que si $S$ est de caractéristique nulle, \deleted{de 3.2 montre que} en utilisant les résultats de Hironaka [\ ] et d'Artin SGA 4 $\overline{\underline{\mathrm{XIX}}}$ \ldots\ . D'autre part, \deleted{\unsure{pour cela,}} lorsqu'on \unsure{disposera|dispose} de la résolution des singularités \deleted{\unsure{sous la forme où}} pour les schémas de type fini sur un corps $k_0$ \inserted{(en supposant \unsure{s. dém. algébrique})}, alors, utilisant le th. de pureté relatif SGA 4 $\overline{\underline{\mathrm{XV}}}$ \ldots\ , on voit \deleted{que le \unsure{même rais.},} \inserted{aisément (en procédant à peu de choses près comme pour 3.5), que le} th. de monodromie est valable \deleted{pour les schémas,} \deleted{\unsure{sous}} sous les formes 3.5 et 3.7, lorsque $S$ est \deleted{local} hensélisé \inserted{\unsure{--|ou} hensélisé strict} d'un schéma de type fini sur \deleted{un corps} $k_0$.

% source lines 24--26
\deleted{\unsure{Par a.}} b) Lorsque $n=1$, \unsure{alors|on a} grâce à Abhyankar [\ ] la condition (ii) de 3.5. est tjs satisfaite. On ne sait apparemment s'il en est de $=$ de la
